% !TeX spellcheck = ru_RU
% !TEX root = vkr.tex

\section*{Введение}
% \thispagestyle{withCompileDate}
% Больше информации можно поискать в слайдах Д.~Кознова~\cite{koznov}.

% Тут  4 части (абзаца) максимум на 2 страницы:
% \begin{enumerate}
% \item Background, known information.
% \item Knowledge gap, unknown information.
% \item  Hypothesis, question, purpose statement.
% \item Approach, plan of attack, proposed solution.
% \begin{itemize}
% \item Последний абзац должен читаться и быть понятным в отрыве от остальных трёх.
% \end{itemize}
% \end{enumerate}

Современные системы связи и радиолокации распространены в разных областях деятельности человека. Речь идет о таких системах как сети высокоскоростного мобильного интернета 4G LTE и 5G или радиостанциях класса DMR, используемых для обеспечения сухопутной мобильной радиосвязи. Неотъемлемой частью таких систем является функция обнаружение сигналов, где под обнаружением понимается определение момента времени начала сигнала и его смещения по частоте относительно центральной частоты приема.

Данная задача является сложной в силу наличия в радиоэфире шумов и помех, неидеальностью приёмной и передающей аппаратуры, а также эффектов многолучевого распространения сигналов (доплеровский сдвиг частоты, искажение формы сигнала).

Оптимальным методом обнаружения сигналов на фоне аддитивного белого Гауссовского шума (АБГШ) является корреляционное обнаружения сигналов \cite{grishin_radio_systems}. Основная идея этого метода заключается в том, что в сигнал на передающей стороне помещается специальная последовательность (преамбула), которая заранее известна как передатчику, так и приемнику. При приеме сигнала осуществляется вычисление корреляции между принятым сигналом и сохраненным в памяти эталонным образцом. В силу того, что корреляция двух сигналов является мерой их сходства, при появлении в принимаемом сигнале преамбулы значение корреляции резко возрастает, что позволяет идентифицировать начало полезного сигнала.

Для вычисления корреляционного интеграла при частоте дискретизации сигнала fs и длине преамбулы в L отсчетов требуется выполнить fs * L операций умножения-сложения (multiply-accumulate, MAC) за секунду. Соответственно, для системы цифровой связи с частотой дискретизации 400 kHz и длиной преамбулы в 1024 отсчета для расчета корреляционного интеграла потребуется вычислитель с производительностью не менее 410 MIPS, что, в принципе, не является проблемой для современного аппаратного обеспечения. 
	
Однако принимаемый сигнал содержит различные искажения, в числе которых особенно важна ошибка по частоте \cite{mengali_sync}. Наличие частотной ошибки в сигнале препятствует расчету корреляции, поскольку преамбула с частотной ошибкой не совпадает с эталонной и, следовательно, полученное значение корреляции будет ниже порогового. Таким образом, наличие частотной ошибки приводит к деградации корреляционного пика и делает обнаружение сигнала невозможным. В связи с этим эффект частотной ошибки должен быть нивелирован, что достигается путем определения величины частотной ошибки и выполнения частотной коррекции сигнала на найденную величину.

Необходимость компенсации частотной ошибки приводит к значительному увеличению вычислительной сложности метода корреляционного обнаружения сигналов. При переборе N различных частотных ошибок необходимо выполнить N раз расчет корреляционного интеграла, что приводит к увеличению объема вычислений в N раз (типовым значением N является 100). Таким образом, реализация метода корреляционного обнаружения сигналов для системы связи, рассмотренной ранее, потребует производительность вычислителя уже в 41000 MIPS, что является вызовом для современных вычислителей.

В качестве основного вычислителя для систем радиолокации и связи могут выступать цифровые сигнальные процессоры, интегральные схемы и ПЛИС \cite{adams_architecture_choice}. Основными производителями сигнальных процессоров являются компании Texas Instruments и Analog Devices (США), которые выпускаю семейства процессоров TMS320C64xx \cite{tms_overview} и TigerSHARC \cite{tsharc_product_brief} соответственно. Использование сигнальных процессоров оправдано при производстве мелкосерийных или средне серийных изделий, требующих для своей работы умеренную производительность. В случае крупносерийного производства экономически оправданным становится переход к использованию интегральных схем. В данном сценарии также появляется интерес к ПЛИС (например AMD 7 series FPGA \cite{xilinx_7_overview} или  Altera Cyclone IV FPGA \cite{cyclone_iv_overview}) как к средству прототипирования целевой интегральной схемы. Также выбор  ПЛИС может быть обусловлен высокой вычислительной сложностью решаемой задачи, для которой производительности сигнального процессора будет недостаточно \cite{fpgas_dsp_applications}.

Одной из платформ, сочетающей интерфейс с радиоканалом (микросхема, обеспечивающая прием и передачу радиосигналов), ПЛИС и процессор, является технология программно-определяемой радиосистемы (Software Defined Radio, SDR) \cite{adi_sdr_solutions}. В рамках данной технологии реализация ряда компонентов системы связи выполняется средствами программными, а не аналоговыми, что позволяет осуществлять перенастройку широкого спектра параметров прямо во время работы SDR. Данным обстоятельством, в совокупности со способностью современных ПЛИС к выполнению частичной переконфигурации без нарушения работоспособности остальной части ПЛИС, объясняется высокая гибкость SDR-решений и интерес к использованию программно-определяемых радиосистем в качестве платформы для разработки и прототипирования систем цифровой связи.

В рамках технологии SDR реализация ресурсоемкого метода корреляционного обнаружения сигналов непосредственно на ПЛИС позволяет разгрузить центральный процессор системы для выполнения других задач, таких как демодуляция, декодирование, протокольная обработка и управление системой в целом.

АО “Концерн ГРАНИТ”  занимается разработкой современных систем связи и радиоэлектроники для профильных министерств и ведомств Российской Федерации. В рамках работ по созданию беспроводных систем связи было необходимо разработать обнаружитель цифровых сигналов, работающий на ПЛИС. В качестве платформы для реализации обнаружителя было выдвинуто требование по использованию устройства Pluto SDR+, которое является клоном ADALM-PLUTO SDR \cite{adalm_pluto_board} производства Analog Devices, и представляет собой доступную реализацию SDR на базе системы на кристалле (System on a Chip, SoC) Xilinx Zynq-7000 \cite{zynq_7000_overview}, сочетающую программируемую логику и двухъядерный процессор ARM.


